\documentclass[12pt]{article}

% ── Packages ───────────────────────────────────────────────────────────────
\usepackage[T1]{fontenc}
\usepackage[utf8]{inputenc}
\usepackage[top=1.1in, bottom=1.1in, left=1.25in, right=1.25in]{geometry}
\usepackage{amsmath, amssymb, amsthm}
\usepackage{mathtools}
\usepackage{bm}
\usepackage{xcolor}
\usepackage{booktabs}
\usepackage{array}
\usepackage{enumitem}
\usepackage{titlesec}
\usepackage{fancyhdr}
\usepackage{hyperref}
\usepackage{cleveref}
\usepackage{algorithm}
\usepackage{algpseudocode}
\usepackage{setspace}
\usepackage{parskip}
\usepackage{tcolorbox}
\usepackage{multirow}
\usepackage{tabularx}
\usepackage{longtable}

\tcbuselibrary{skins, breakable, theorems}

% ── Colors ─────────────────────────────────────────────────────────────────
\definecolor{accent}{HTML}{2D1B4E}
\definecolor{accent2}{HTML}{6B3FA0}
\definecolor{lightbg}{HTML}{F3EFF8}
\definecolor{rulecol}{HTML}{C8B8E8}
\definecolor{notecol}{HTML}{1E5C3A}
\definecolor{notebg}{HTML}{EAF4ED}
\definecolor{warnbg}{HTML}{FDF3E7}
\definecolor{warncol}{HTML}{7A4100}
\definecolor{eqbg}{HTML}{F3EFF8}
\definecolor{graytext}{HTML}{444444}

% ── Hyperref ───────────────────────────────────────────────────────────────
\hypersetup{
  colorlinks=true,
  linkcolor=accent2,
  citecolor=accent,
  urlcolor=accent2,
  pdftitle={EVE: Evolutionary Virtual Exploration -- Research Proposal},
  pdfauthor={Research Proposal}
}

% ── Section style ──────────────────────────────────────────────────────────
\titleformat{\section}
  {\color{accent}\large\bfseries\vspace{4pt}}
  {\thesection.}{0.6em}{}[{\color{rulecol}\vspace{2pt}\hrule height 0.6pt\vspace{2pt}}]

\titleformat{\subsection}
  {\color{accent2}\normalsize\bfseries}
  {\thesubsection.}{0.5em}{}

\titleformat{\subsubsection}
  {\color{graytext}\normalsize\bfseries\itshape}
  {\thesubsubsection.}{0.5em}{}

\titlespacing{\section}{0pt}{18pt}{8pt}
\titlespacing{\subsection}{0pt}{12pt}{4pt}
\titlespacing{\subsubsection}{0pt}{8pt}{3pt}

% ── Header / Footer ────────────────────────────────────────────────────────
\pagestyle{fancy}
\fancyhf{}
\renewcommand{\headrulewidth}{0.4pt}
\fancyhead[L]{\small\color{graytext}\textit{EVE} \textemdash\
              Evolutionary Virtual Exploration}
\fancyhead[R]{\small\color{graytext}Research Proposal}
\fancyfoot[C]{\small\color{graytext}\thepage}

% ── tcolorbox styles ───────────────────────────────────────────────────────
\tcbset{
  eqstyle/.style={
    colback=eqbg, colframe=accent2, arc=2pt,
    boxrule=0.5pt, left=6pt, right=6pt, top=5pt, bottom=5pt,
    fonttitle=\footnotesize\bfseries\color{accent2}, attach title to upper={\\[4pt]}
  },
  notestyle/.style={
    colback=notebg, colframe=notecol, arc=2pt,
    boxrule=0.5pt, left=6pt, right=6pt, top=4pt, bottom=4pt,
    fontupper=\small\itshape\color{notecol}
  },
  warnstyle/.style={
    colback=warnbg, colframe=warncol, arc=2pt,
    boxrule=0.5pt, left=6pt, right=6pt, top=4pt, bottom=4pt,
    fontupper=\small\bfseries\color{warncol}
  }
}

\newcommand{\eqbox}[2]{%
  \begin{tcolorbox}[eqstyle, title={#1}]
    #2
  \end{tcolorbox}
}

\newcommand{\notebox}[1]{%
  \begin{tcolorbox}[notestyle]
    #1
  \end{tcolorbox}
}

\newcommand{\warnbox}[1]{%
  \begin{tcolorbox}[warnstyle]
    #1
  \end{tcolorbox}
}

% ── Theorem environments ───────────────────────────────────────────────────
\newtheorem{definition}{Definition}
\newtheorem{proposition}{Proposition}
\newtheorem{theorem}{Theorem}
\newtheorem{corollary}{Corollary}
\newtheorem{lemma}{Lemma}
\newtheorem{remark}{Remark}

% ── Macros ─────────────────────────────────────────────────────────────────
\newcommand{\R}{\mathbb{R}}
\newcommand{\E}{\mathbb{E}}
\newcommand{\norm}[1]{\left\lVert#1\right\rVert}
\newcommand{\abs}[1]{\left\lvert#1\right\rvert}
\newcommand{\ind}[1]{\mathbf{1}\!\left[#1\right]}
\newcommand{\Loss}{\mathcal{L}}
\newcommand{\softmax}{\mathrm{softmax}}
\newcommand{\sign}{\mathrm{sign}}
\newcommand{\diag}{\mathrm{diag}}
% ── Spacing helpers ────────────────────────────────────────────────────────
\setlength{\parskip}{6pt}
\setlength{\parindent}{0pt}

% ══════════════════════════════════════════════════════════════════════════
\begin{document}

% ── Title page ─────────────────────────────────────────────────────────────
\begin{center}
  {\color{accent}\LARGE\bfseries EVE}\\[6pt]
  {\color{accent2}\Large Evolutionary Virtual Exploration}\\[8pt]
  {\large Adaptive Optimization via Offspring Direction\\[2pt]
  Selection with Probe-Based Fitness Evaluation}\\[10pt]
  {\normalsize\color{graytext}\bfseries Research Proposal}\\[6pt]
  {\normalsize\color{graytext}\itshape March 2026}\\[6pt]
\end{center}

\noindent{\color{rulecol}\rule{\linewidth}{1.2pt}}

% ── Abstract ───────────────────────────────────────────────────────────────
\begin{center}
  {\bfseries Abstract}
\end{center}
\noindent
Adam and its variants dominate deep learning optimization by maintaining
per-parameter adaptive learning rates through exponential moving averages
of gradient moments.
Yet Adam commits to a single update direction at every step, cannot detect
when that direction is locally suboptimal, and structurally neglects
parameter dimensions where gradient signal is noisy.
We propose \textbf{EVE}---Evolutionary Virtual Exploration---an optimizer
that embeds evolutionary selection pressure directly inside the parameter
update rule.
At each step, EVE constructs $K$ \emph{offspring directions}: competing
update strategies derived analytically from the optimizer's existing
moment estimates.
A probe-based fitness evaluation tests each offspring on the training batch
via $K$ forward-only passes, and soft natural selection via
temperature-scaled softmax combines the offspring into a single update
weighted by empirically measured fitness.
An evolutionary memory---the \emph{strength signal}---tracks per-dimension
historical selection success and drives the construction of complementary
offspring that target neglected dimensions.
The framework is parameterized by a single integer $K$ (brood size):
at $K=1$, EVE \emph{exactly} degenerates to AdamW with zero overhead;
at $K>1$, EVE adaptively selects among $K$ strategies at the cost of
$K$ additional forward passes on the training batch.

The offspring direction set is informed by a pilot study on CIFAR-10
with ResNet-18, which identified two critical constraints for offspring
design: (i)~momentum smoothing is non-negotiable for training
stability, and (ii)~each offspring must occupy a distinct strategic
niche to prevent any single direction from bearing disproportionate
selection load.
EVE's four canonical offspring accordingly cover stable exploitation
(AdamW direction), trajectory smoothing (slow momentum), targeted
dimension recovery (interpolated complementary), and sharpness
awareness (virtual SAM perturbation).
EVE preserves Adam's convergence guarantees at $K=1$ and provides
exploration--exploitation control via adaptive selection
temperature, with no additional hyperparameters beyond the brood size.
This proposal presents the theoretical framework, algorithmic design,
pilot results, and analysis of expected behavior for EVE, and outlines
the experimental validation required to establish its empirical
effectiveness.

\vspace{4pt}
\noindent{\color{rulecol}\rule{\linewidth}{0.6pt}}

% ══════════════════════════════════════════════════════════════════════════
\section{Introduction and Motivation}
\label{sec:intro}

The Adam optimizer~\cite{kingma2015adam} and its decoupled weight decay
variant AdamW~\cite{loshchilov2019decoupled} have become the default
optimizers for training deep neural networks.
Their success relies on two exponential moving averages (EMAs): a first
moment $m_t$ (momentum) and a second moment $v_t$ (squared gradient
magnitude), which together produce a per-parameter adaptive learning rate.
This mechanism is simple, memory-efficient, and broadly effective.

Yet Adam has structural limitations that no amount of learning rate
tuning can fully address.

\textbf{Single-direction commitment.}
At every step, Adam computes exactly one update direction:
$-\hat{m}_t / (\sqrt{\hat{v}_t} + \varepsilon)$.
If this direction is locally suboptimal---sliding along a ravine
instead of cutting across it, or stepping toward a sharp minimum
instead of escaping it---Adam has no internal mechanism to detect the
error.
The optimizer follows its single signal blindly, relying entirely on
future gradients to correct course.

\textbf{Sharp minima entrapment.}
In regions of high curvature, both the gradient magnitude and $v_t$
are large.
Adam's normalization $\hat{m}_t / \sqrt{\hat{v}_t}$ cancels the
sharpness signal: the optimizer takes small, careful steps precisely
when it should be leaping out.
This structural blindness to loss landscape geometry is the failure mode
that Sharpness-Aware Minimization (SAM)~\cite{foret2021sharpnessaware}
was designed to address, at the cost of doubled forward--backward passes.

\textbf{Dimension neglect.}
On parameter dimensions where the gradient oscillates in sign, $m_t$
averages toward zero while $v_t$ remains large.
The effective learning rate
$\eta \cdot |\hat{m}_{t,d}| / (\sqrt{\hat{v}_{t,d}} + \varepsilon)$
collapses, and the dimension becomes effectively frozen.
With $\beta_2 = 0.999$, the half-life of $v_t$ is approximately
693 steps; a neglected dimension can remain stalled for hundreds of
updates.

\textbf{Hyperparameter fragility.}
Adam's performance is highly sensitive to the learning rate $\eta$.
The optimal value can vary by an order of magnitude between tasks,
and practitioners routinely spend significant compute on learning rate
sweeps, warmup schedules, and cosine annealing---all patches for an
optimizer that cannot self-correct.

\medskip

These are not bugs in Adam's implementation; they are structural
consequences of committing to a single update formula per step.
We propose that EVE can address these limitations by introducing a minimal evolutionary layer:
multiple candidate directions are spawned, tested, and selected at
every step, so the optimizer can adapt its \emph{strategy}---not just
its step size---to the local loss landscape.

The key insight is that evolutionary selection can be embedded directly
in the update rule without maintaining a population of models.
EVE spawns \emph{offspring directions}---virtual candidate updates
derived analytically from existing moment buffers---and evaluates their
fitness via $K$ forward-only passes on the training batch.
Each pass perturbs the parameters in-place along one offspring direction,
measures the resulting loss, then restores the original parameters before
the next probe.
Soft selection via temperature-scaled softmax produces a combined
update that is a convex combination of the offspring directions,
and the brood size $K$ provides
a clean, user-facing knob trading compute for adaptivity.

At $K = 1$, only the Adam direction exists, selection is trivial, and
EVE is exactly AdamW with zero overhead.
At $K > 1$, EVE adaptively selects among $K$ strategies at the cost
of $K$ additional forward passes.
No populations, no evolutionary operators applied to model weights,
no surrogate maps---just $K$ directions, $K$ probes, and a softmax.


% ══════════════════════════════════════════════════════════════════════════
\section{Pilot Study: Offspring Dominance on CIFAR-10}
\label{sec:pilot}

To inform the offspring direction design, we conducted a pilot study
training ResNet-18 on CIFAR-10 with $K = 4$ offspring.
This early experiment tested a naive offspring set in which
$\mathbf{d}_2$ was the raw gradient direction
($-g_t / (\sqrt{\hat{v}_t} + \varepsilon)$, no momentum smoothing),
$\mathbf{d}_3$ was a raw-gradient complementary direction
($-g_t \odot (1 - s_t) / (\sqrt{\hat{v}_t} + \varepsilon)$), and
$\mathbf{d}_4$ was a curvature-proportional contrarian
($-\sign(g_t) \odot \sqrt{\hat{v}_t} / (\max(\sqrt{\hat{v}_t}) +
\varepsilon)$).
The dominance summary:

\begin{center}
\small
\begin{tabular}{@{}llrrr@{}}
  \toprule
  Offspring & Role & Dominant (\%) & Avg Weight & Avg Fitness \\
  \midrule
  $\mathbf{d}_1$ (Adam)          & Momentum-smoothed step     & 25.2 & 0.2944 &  0.0941 \\
  $\mathbf{d}_2$ (Raw Gradient)  & Fresh signal, no momentum  &  0.0 & 0.0262 & $-1.0149$ \\
  $\mathbf{d}_3$ (Complementary) & Targets weak dimensions    & 66.2 & 0.4859 &  0.1232 \\
  $\mathbf{d}_4$ (Contrarian)    & High-curvature escape      &  8.7 & 0.1936 &  0.0480 \\
  \bottomrule
\end{tabular}
\end{center}

Three observations shaped the final offspring design:

\textbf{Observation 1: Momentum is non-negotiable.}
The raw gradient offspring ($\mathbf{d}_2$) achieved fitness of
$-1.015$---not marginally bad but catastrophically harmful,
\emph{increasing} the loss when applied.
On CIFAR-10 with ResNet-18 and BatchNorm, the per-step gradient is
too noisy for direct use as an update direction.
Any offspring must incorporate momentum smoothing.

\textbf{Observation 2: The complementary offspring was overloaded.}
$\mathbf{d}_3$ dominated 66.2\% of steps, carrying more than double
the load of the next most dominant offspring.
This indicates that Adam's momentum systematically stalls on
a majority of parameter dimensions---the $m_t \to 0$ problem on
sign-oscillating gradients---and the complementary direction was the
only offspring capable of recovering these dimensions.
However, because $\mathbf{d}_3$ used the raw gradient $g_t$ (inheriting
$\mathbf{d}_2$'s noise problem on the dimensions where it was most
active), it was dominant partly by default: it was the sole non-Adam
strategy with positive fitness.

\textbf{Observation 3: The contrarian had limited value.}
$\mathbf{d}_4$ was dominant only 8.7\% of the time with the lowest
positive fitness (0.048).
Its curvature-proportional design was structurally different from the
other offspring, but the empirical signal suggests its slot could be
used more effectively.

\medskip

\warnbox{%
\textbf{Design principle.}
The offspring set should contain no dead directions (zero dominance)
and no overloaded compensators ($>$50\% dominance).
Each offspring must occupy a distinct strategic niche and incorporate
sufficient smoothing to survive the noise profile of real training.
}


% ══════════════════════════════════════════════════════════════════════════
\section{Related Work}
\label{sec:related}

\subsection{Adaptive Gradient Methods}

Adam~\cite{kingma2015adam} maintains per-parameter first and second
moment estimates; AdamW~\cite{loshchilov2019decoupled} decouples weight
decay from the adaptive update.
RAdam~\cite{liu2020radam} addresses Adam's variance issue during early
training via a rectified update.
AdaFactor~\cite{shazeer2018adafactor} reduces memory by factoring the
second-moment matrix.
Lion~\cite{chen2024lion} uses only the sign of the momentum, eliminating
the second moment entirely.
All commit to a single update formula per step.

\subsection{Sharpness-Aware Optimization}

SAM~\cite{foret2021sharpnessaware} seeks parameters where the loss is
uniformly low in a neighborhood, requiring two forward--backward passes
per step (100\% overhead).
GSAM~\cite{zhuang2022surrogate} and LookSAM~\cite{liu2022towards}
reduce this cost via periodic perturbation or surrogate gradients.
EVE's $\mathbf{d}_4$ (virtual SAM offspring) provides a softer form of
sharpness awareness: the perturbed direction is evaluated via a
forward-only probe on the training batch, and downweighted by soft
selection if the perturbation reveals sharpness.

\subsection{Learning Rate Adaptation}

D-Adaptation~\cite{defazio2023dadaptation} and Prodigy~\cite{mishchenko2024prodigy}
automatically estimate the optimal learning rate during training.
Schedule-Free optimizers~\cite{defazio2024schedulefree} eliminate the
need for learning rate schedules.
EVE is orthogonal to these: it adapts the \emph{direction}, not just
the magnitude, and can be composed with any learning rate schedule or
automatic rate estimator.

\subsection{Multi-Candidate and Ensemble Methods}

Lookahead~\cite{zhang2019lookahead} maintains slow and fast weights,
effectively exploring two update strategies.
EVE's $\mathbf{d}_2$ (slow-momentum offspring) internalizes the
Lookahead principle as a competing direction within the selection
mechanism, rather than requiring an external wrapper optimizer.
Stochastic Weight Averaging (SWA)~\cite{izmailov2018averaging} averages
multiple points along the trajectory.
Gradient Disagreement~\cite{mansilla2021domain} uses per-sample gradient
variance to detect distribution shift.
EVE generalizes the multi-candidate idea: $K$ strategies compete at
every step, selected by empirical fitness rather than heuristic averaging.

\subsection{Evolutionary Optimization for Neural Networks}

Evolution Strategies (ES)~\cite{salimans2017evolution} demonstrated that
population-based search can match gradient methods on RL benchmarks.
CMA-ES~\cite{hansen2001cma} adapts a full covariance mutation matrix
but scales as $\mathcal{O}(D^2)$.
EVE distills the evolutionary insight---multiple strategies, fitness-based
selection, complementarity-driven exploration---into a single-model
optimizer with zero population overhead.

\subsection{Symbolic and Learned Optimizer Discovery}

The learning-to-optimize (L2O) paradigm~\cite{andrychowicz2016l2l}
parameterizes the update rule as a neural network trained via
meta-learning over a distribution of tasks.
VeLO~\cite{metz2022velo} scaled this approach using approximately
4,000 TPU-months of compute, though independent evaluation raised
concerns about generalization~\cite{kaddour2023velo}.
Symbolic search methods~\cite{bello2017neural, chen2024lion} discover
interpretable update rules but produce \emph{fixed} formulas that
inherit the single-direction limitation.

EVE takes a fundamentally different approach.
Rather than searching for a better fixed update rule at enormous
computational cost, EVE provides a lightweight runtime mechanism
for adaptively selecting among multiple update strategies at every
step.
Discovered optimizers (such as Lion) could themselves serve as
offspring directions within EVE's selection mechanism.


% ══════════════════════════════════════════════════════════════════════════
\section{Notation and Preliminaries}
\label{sec:notation}

Let $\theta_t \in \R^D$ denote the parameter vector at step $t$,
$\Loss(\theta; B)$ the loss evaluated on mini-batch $B$, and
$g_t = \nabla_\theta \Loss(\theta_t; B_t)$ the stochastic gradient.

Adam maintains two EMA buffers:
\begin{align}
  m_t &= \beta_1 m_{t-1} + (1 - \beta_1)\, g_t
  \label{eq:adam_m} \\
  v_t &= \beta_2 v_{t-1} + (1 - \beta_2)\, g_t^2
  \label{eq:adam_v}
\end{align}
with bias-corrected estimates
$\hat{m}_t = m_t / (1 - \beta_1^t)$ and
$\hat{v}_t = v_t / (1 - \beta_2^t)$.
All operations are element-wise over dimensions $d \in \{1, \ldots, D\}$.

EVE introduces an additional slow momentum buffer:
\begin{equation}
  m_{t,\mathrm{slow}} = \beta_{1,\mathrm{slow}}\, m_{t-1,\mathrm{slow}}
    + (1 - \beta_{1,\mathrm{slow}})\, g_t
  \label{eq:m_slow}
\end{equation}
with $\beta_{1,\mathrm{slow}} = 0.999$ (default) and bias-corrected
estimate $\hat{m}_{t,\mathrm{slow}} = m_{t,\mathrm{slow}} /
(1 - \beta_{1,\mathrm{slow}}^t)$.
This buffer tracks the low-frequency trend of the gradient signal,
producing a trajectory-smoothed direction that is permanently
differentiated from the standard momentum ($\beta_1 = 0.9$).

The AdamW update is:
\begin{equation}
  \theta_{t+1} = \theta_t
    - \eta\,\frac{\hat{m}_t}{\sqrt{\hat{v}_t} + \varepsilon}
    - \eta\lambda\,\theta_t
  \label{eq:adamw}
\end{equation}
where $\eta$ is the learning rate, $\lambda$ the weight decay, and
$\varepsilon$ a numerical stabilizer.


% ══════════════════════════════════════════════════════════════════════════
\section{Proposed Method: The EVE Framework}
\label{sec:method}

EVE augments the standard adaptive gradient update with three mechanisms:
(1)~offspring direction construction, (2)~probe-based fitness evaluation,
and (3)~soft natural selection.
A fourth mechanism, the strength signal, provides evolutionary memory
that links generations across steps.

\subsection{Offspring Direction Construction}
\label{ssec:offspring}

At each step $t$, after computing $g_t$, $\hat{m}_t$, $\hat{v}_t$,
$\hat{m}_{t,\mathrm{slow}}$, and the strength signal $s_t$
(\cref{ssec:strength}), EVE constructs $K$ offspring directions.
Each offspring represents a distinct optimization strategy, derived
analytically from the existing moment buffers at negligible computational
cost.

The offspring set is guided by three design constraints derived from
the pilot study (\cref{sec:pilot}):
(i)~every offspring must incorporate momentum smoothing (raw gradient
directions are catastrophically unfit under typical training noise),
(ii)~offspring roles must be non-overlapping so that no single
direction bears disproportionate selection load, and
(iii)~every offspring slot must contain a direction that is
structurally viable (positive expected fitness).

\eqbox{Definition 1 --- Offspring Directions}{
Let $K \geq 1$ be the brood size.
The offspring directions
$\mathbf{d}_1, \ldots, \mathbf{d}_K \in \R^D$ are defined as:
\begin{align}
  \mathbf{d}_1 &= -\,\frac{\hat{m}_t}{\sqrt{\hat{v}_t} + \varepsilon}
  &&\text{(Baseline --- AdamW direction)}
  \label{eq:offspring_baseline} \\[8pt]
  \mathbf{d}_2 &= -\,\frac{\hat{m}_{t,\mathrm{slow}}}
    {\sqrt{\hat{v}_t} + \varepsilon}
  &&\text{(Slow Momentum --- trajectory smoother)}
  \label{eq:offspring_slow} \\[8pt]
  \mathbf{d}_3 &= -\,\frac{
    \bigl[\alpha\,\hat{m}_t + (1 - \alpha)\,g_t\bigr]
    \odot (1 - s_t)
  }{\sqrt{\hat{v}_t} + \varepsilon}
  &&\text{(Interpolated Complementary --- anti-stall)}
  \label{eq:offspring_complementary} \\[8pt]
  \mathbf{d}_4 &= \mathbf{d}_1
    + \gamma \cdot \sign(g_t)
  &&\text{(Virtual SAM --- sharpness probe)}
  \label{eq:offspring_sam}
\end{align}
where $\odot$ denotes element-wise multiplication,
$s_t \in [0, 1]^D$ is the strength signal (\cref{ssec:strength}),
$\alpha \in (0, 1)$ is the momentum interpolation coefficient
(default: $0.5$), and $\gamma > 0$ is the SAM perturbation scale
(default: $0.01 \cdot \eta$).
}

\notebox{%
\textbf{Design rationale.}
\begin{itemize}[leftmargin=1.6em, itemsep=4pt]
  \item \textbf{$\mathbf{d}_1$ (Baseline):} The standard
    momentum-smoothed adaptive step.
    It guarantees that at $K = 1$, EVE safely degenerates to AdamW
    with zero overhead.
    It acts as the ``stable exploiter'' of the brood.
  \item \textbf{$\mathbf{d}_2$ (Slow Momentum):}
    Uses a secondary momentum buffer $\hat{m}_{t,\mathrm{slow}}$
    computed with a much higher decay rate
    ($\beta_{1,\mathrm{slow}} = 0.999$ vs.\ $\beta_1 = 0.9$).
    This tracks the trajectory's low-frequency trend while
    $\mathbf{d}_1$ tracks the high-frequency signal.
    When $\mathbf{d}_1$ oscillates between valley walls, the probe
    naturally selects $\mathbf{d}_2$ to pull the trajectory back
    to the valley center---internalizing the Lookahead
    principle~\cite{zhang2019lookahead} within EVE's selection
    mechanism.
    The pilot study (\cref{sec:pilot}) demonstrated that any offspring
    lacking momentum smoothing achieves catastrophically negative
    fitness; the slow momentum design ensures $\mathbf{d}_2$ is
    always a viable competitor.
  \item \textbf{$\mathbf{d}_3$ (Interpolated Complementary):}
    Targets dimensions where Adam's momentum has stalled
    ($s_t \approx 0$) by injecting a blend of momentum and fresh
    gradient: $\alpha\,\hat{m}_t + (1 - \alpha)\,g_t$ with
    $\alpha = 0.5$.
    The momentum interpolation dampens sub-batch noise while
    preserving the ability to break out of momentum-stalled
    dimensions.
    The pilot study showed that a pure-gradient complementary
    direction dominated 66\% of steps---effective but overloaded.
    The momentum blend distributes the anti-stall workload
    more evenly across the brood.
  \item \textbf{$\mathbf{d}_4$ (Virtual SAM):}
    Steps in the Adam direction plus a small sign perturbation
    proportional to $\gamma$.
    True SAM~\cite{foret2021sharpnessaware} requires a costly
    second forward--backward pass to evaluate a worst-case
    neighborhood perturbation.
    EVE achieves this ``virtually'': if the local region is flat,
    $\mathbf{d}_4 \approx \mathbf{d}_1$ and the perturbation is
    harmless.
    If the region is sharp, the perturbation causes poor probe
    fitness, $\mathbf{d}_4$ is downweighted, and the combined
    update is biased away from the sharp direction---achieving
    sharpness awareness through EVE's existing probe mechanism
    at no additional backward-pass cost.
\end{itemize}
}

\subsubsection{Role Differentiation}

A key design goal of the offspring set is that each offspring
occupies a distinct niche in strategy space, so that the selection
mechanism has meaningful alternatives to choose from at every step.

\begin{center}
\small
\begin{tabular}{@{}lll@{}}
  \toprule
  Offspring & Strategy Niche & Wins When \ldots \\
  \midrule
  $\mathbf{d}_1$ (Baseline)
    & Standard adaptive momentum
    & Landscape is well-conditioned; Adam works \\
  $\mathbf{d}_2$ (Slow Momentum)
    & Low-frequency trajectory tracking
    & $\mathbf{d}_1$ oscillates between valley walls \\
  $\mathbf{d}_3$ (Complementary)
    & Targeted dimension recovery
    & Momentum has stalled on sign-oscillating dims \\
  $\mathbf{d}_4$ (Virtual SAM)
    & Sharpness-aware perturbation
    & Current minimum is wide and generalizable \\
  \bottomrule
\end{tabular}
\end{center}

\notebox{%
\textbf{Offspring 1 is always the Adam direction.}
Each subsequent offspring adds a progressively more exploratory strategy,
and all share the curvature denominator $\sqrt{\hat{v}_t} + \varepsilon$
so that selection compares directional quality rather than magnitude
artifacts.
Increasing $K$ monotonically increases exploration without altering the
baseline behavior; at $K = 1$, only the Adam direction exists.
For $K > 4$, additional offspring can be built by using sign-momentum
(analogous to Lion~\cite{chen2024lion}), applying layer-aware strategies,
or adding stochastic perturbations around $\mathbf{d}_1$.
}


\subsection{Probe-Based Fitness Evaluation}
\label{ssec:probe}

The core of EVE's selection mechanism is a \emph{probe}: a forward-only
evaluation of each offspring direction on the training batch, producing
a fitness ranking.

\subsubsection{Control-Variate Fitness}

Before probing, a single forward pass at the current parameters
$\theta_t$ establishes a baseline loss on $B_t$:

\eqbox{Definition 2 --- Probe Fitness}{
\begin{equation}
  \mathcal{F}_k = \Loss(\theta_t;\, B_t)
    \;-\; \Loss(\theta_t + \eta\,\mathbf{d}_k;\, B_t)
  \label{eq:fitness}
\end{equation}
\smallskip
Positive $\mathcal{F}_k$ indicates that offspring $k$ reduces loss.
The shared baseline $\Loss(\theta_t; B_t)$ acts as a \emph{control variate}:
batch noise affects all offspring identically and cancels in the
relative ranking.
}

\begin{proposition}[Noise cancellation]
\label{prop:noise}
Let $\xi$ denote the additive batch noise:
$\Loss(\theta; B_t) = \Loss(\theta; \mathcal{D}) + \xi(\theta, B_t)$.
The fitness difference between offspring $k$ and $j$ is:
\begin{equation*}
  \mathcal{F}_k - \mathcal{F}_j
  = \bigl[\Loss(\theta_t + \eta\,\mathbf{d}_j; B_t)
    - \Loss(\theta_t + \eta\,\mathbf{d}_k; B_t)\bigr]
\end{equation*}
The baseline noise $\xi(\theta_t, B_t)$ cancels exactly.
The residual noise $\xi(\theta_t + \eta\mathbf{d}_k, B_t)
- \xi(\theta_t + \eta\mathbf{d}_j, B_t)$ is
$\mathcal{O}(\eta \cdot \norm{\mathbf{d}_k - \mathbf{d}_j} \cdot
\sigma_\xi / \sqrt{|B_t|})$, which is second-order in the learning rate
and negligible for typical $\eta \leq 10^{-3}$.
\end{proposition}

\subsubsection{Probe Loop}

Each offspring is evaluated sequentially via an in-place
perturb--forward--restore loop:

\begin{enumerate}[leftmargin=1.8em, itemsep=3pt]
  \item \textbf{Perturb.} Apply the candidate step in-place:
    $\theta \leftarrow \theta_t + \eta\,\mathbf{d}_k$.
  \item \textbf{Forward.} Run a forward-only pass on $B_t$;
    record $L_k = \Loss(\theta_t + \eta\,\mathbf{d}_k;\, B_t)$.
  \item \textbf{Restore.} Copy the saved parameters back:
    $\theta \leftarrow \theta_t$.
\end{enumerate}

This approach avoids maintaining $K$ copies of the parameters in memory
(only one backup copy is needed), and works correctly with layers that
maintain running statistics (e.g.\ BatchNorm), which can fail under
functional-call parallelization.
No backward pass is performed during the probe.

\eqbox{Proposition 1 --- Computational Cost}{
The probe evaluation adds one baseline forward pass plus $K$ probe
forward passes, all on the full batch $B_t$:
\begin{equation}
  C_{\mathrm{probe}} = (1 + K)\,C_{\mathrm{fwd}}^{\mathrm{full}}
  \label{eq:probe_cost}
\end{equation}
\smallskip
Since a standard training step costs
$C_{\mathrm{fwd}}^{\mathrm{full}} + C_{\mathrm{bwd}}^{\mathrm{full}}
\approx 3\,C_{\mathrm{fwd}}^{\mathrm{full}}$,
the overhead relative to AdamW is approximately
$\frac{1+K}{3} \times 100\%$, or roughly $\sim$167\% for $K = 4$.
No backward pass is required for the probe.
}


\subsection{Soft Natural Selection}
\label{ssec:selection}

The fitness values $\{\mathcal{F}_k\}_{k=1}^K$ are converted to
offspring weights via temperature-scaled softmax:

\eqbox{Definition 3 --- Offspring Selection Weights}{
\begin{equation}
  w_k = \frac{
    \exp\!\bigl(\beta_{\mathrm{sel}}\,\mathcal{F}_k\bigr)
  }{
    \sum_{j=1}^{K}
    \exp\!\bigl(\beta_{\mathrm{sel}}\,\mathcal{F}_j\bigr)
  }
  \label{eq:selection_weights}
\end{equation}
\smallskip
where $\beta_{\mathrm{sel}} > 0$ is the \emph{selection temperature}.
High $\beta_{\mathrm{sel}}$ concentrates weight on the fittest offspring
(exploitation); low $\beta_{\mathrm{sel}}$ distributes weight uniformly
(exploration).
}

\subsubsection{$K = 1$ Degeneracy}

\begin{proposition}[$K = 1$ recovers AdamW]
\label{prop:k1}
When $K = 1$, the only offspring is $\mathbf{d}_1$ (the AdamW direction).
The softmax of a single element yields $w_1 = 1$.
The probe evaluation is skipped (single candidate requires no ranking).
The update reduces to:
\begin{equation}
  \theta_{t+1} = \theta_t + \eta\,\mathbf{d}_1 - \eta\lambda\,\theta_t
  = \theta_t - \eta\,\frac{\hat{m}_t}{\sqrt{\hat{v}_t} + \varepsilon}
    - \eta\lambda\,\theta_t
  \label{eq:k1_adam}
\end{equation}
which is exactly the AdamW update~(\cref{eq:adamw}).
The slow momentum buffer $m_{t,\mathrm{slow}}$ is still updated
(to maintain readiness for switching to $K > 1$) but is not used
in the parameter update.
\end{proposition}

\subsubsection{The EVE Update}

The selected update direction is the fitness-weighted combination:

\eqbox{Definition 4 --- The EVE Update}{
\begin{equation}
  \theta_{t+1} = \theta_t
    + \eta \sum_{k=1}^{K} w_k\,\mathbf{d}_k
    - \eta\lambda\,\theta_t
  \label{eq:eve_update}
\end{equation}
\smallskip
This is a convex combination of offspring directions (since
$w_k \geq 0$, $\sum_k w_k = 1$) scaled by the learning rate, plus
decoupled weight decay.
When all $\mathcal{F}_k$ are equal (no offspring is better),
$w_k = 1/K$ and the update averages the strategies---a principled
fallback that diversifies the step direction.
}


\subsection{Strength Signal: Evolutionary Memory}
\label{ssec:strength}

The strength signal $s_t \in [0, 1]^D$ tracks per-dimension historical
selection success, providing evolutionary memory that links generations
(steps) across time.

\eqbox{Definition 5 --- Strength Signal Update}{
At step $t+1$, after computing $g_{t+1}$, the strength signal is updated:
\begin{equation}
  s_{t+1,d} = \gamma_s\, s_{t,d}
    + (1 - \gamma_s)\,\sigma\!\bigl(
      \delta_t \cdot \sign(\Delta\theta_{t,d})
      \cdot \sign(-g_{t+1,d})
    \bigr)
  \label{eq:strength}
\end{equation}
\smallskip
where:
\begin{itemize}[leftmargin=1.6em, itemsep=2pt]
  \item $\delta_t = \Loss(\theta_t; B_t) - \Loss(\theta_{t+1}; B_{t+1})$
    is the inter-step loss improvement (positive = loss decreased),
  \item $\Delta\theta_{t,d} = \theta_{t+1,d} - \theta_{t,d}$ is the
    actual parameter displacement on dimension $d$,
  \item $\sigma(\cdot)$ is the sigmoid function,
  \item $\gamma_s \in [0, 1)$ is the strength decay rate
    (default: $0.99$).
\end{itemize}
Initialize $s_{0,d} = 0.5$ (neutral prior).
}

\textbf{Interpretation.}
The indicator
$\sign(\Delta\theta_{t,d}) \cdot \sign(-g_{t+1,d})$ fires positively
when the update direction on dimension $d$ agreed with the next step's
negative gradient---i.e., the optimizer moved in a direction that the
next gradient confirms was useful.
Multiplying by $\delta_t$ scales this by overall step quality: a
dimension gets credit for agreeing with the descent direction only
when the step actually reduced loss.
Through the sigmoid, $s_{t,d} \to 1$ on dimensions where updates
consistently help, and $s_{t,d} \to 0$ on dimensions where the
optimizer is failing.

\textbf{Role in offspring construction.}
The interpolated complementary offspring (\cref{eq:offspring_complementary})
multiplies the blended signal by $(1 - s_t)$: dimensions with
$s_t \approx 1$ (Adam is already succeeding) are suppressed, while
dimensions with $s_t \approx 0$ (Adam is failing) are amplified.
This creates targeted exploration on exactly the dimensions that need it,
while the momentum interpolation ($\alpha = 0.5$) ensures the injected
signal retains trajectory stability.

\notebox{%
\textbf{Evolutionary analogy.}
The strength signal is analogous to \emph{fitness inheritance} in
biological evolution: offspring inherit information about which traits
(dimensions) contributed to parental survival, guiding the next
generation's variation toward productive dimensions.
EVE encodes this as a continuous signal derived from natural training
dynamics, requiring no explicit survival indicators or population-level
statistics.
}


\subsection{Adaptive Selection Temperature}
\label{ssec:temperature}

The selection temperature $\beta_{\mathrm{sel}}$ controls the
exploration--exploitation balance and adapts automatically based on
selection entropy.

\eqbox{Definition 6 --- Adaptive Selection Temperature}{
Let $H_t = -\sum_{k=1}^K w_k \log w_k$ be the selection entropy
at step $t$, and $H_{\max} = \log K$ the maximum entropy (uniform
selection).
The target entropy is $H^* = \rho \cdot H_{\max}$ for
$\rho \in (0, 1)$ (default: $\rho = 0.5$).

The temperature adapts via multiplicative update:
\begin{equation}
  \beta_{\mathrm{sel}}^{(t+1)} = \beta_{\mathrm{sel}}^{(t)} \cdot
    \exp\!\bigl(\alpha_\beta\,(H_t - H^*)\bigr)
  \label{eq:temperature_adapt}
\end{equation}
\smallskip
where $\alpha_\beta > 0$ is the adaptation rate (default: $0.01$).
When entropy is too high ($H_t > H^*$: selection too uniform),
$\beta_{\mathrm{sel}}$ increases to sharpen selection.
When entropy is too low ($H_t < H^*$: one offspring dominates),
$\beta_{\mathrm{sel}}$ decreases to encourage exploration.
Clamp $\beta_{\mathrm{sel}} \in [\beta_{\min}, \beta_{\max}]$ for
stability (defaults: $[0.1,\, 100]$).
}


% ══════════════════════════════════════════════════════════════════════════
\section{Theoretical Properties}
\label{sec:theory}

\subsection{Dominance Guarantee}

\begin{theorem}[EVE dominates any single offspring]
\label{thm:dominance}
Let $\mathcal{F}_k$ be the true (population-level) fitness of offspring
$k$.
The expected fitness of the EVE update is:
\begin{equation}
  \E\bigl[\mathcal{F}_{\mathrm{EVE}}\bigr]
  = \sum_{k=1}^K w_k\,\mathcal{F}_k
  \geq \min_{k} \mathcal{F}_k
  \label{eq:dominance}
\end{equation}
with equality only when all offspring have equal fitness.
Furthermore, as $\beta_{\mathrm{sel}} \to \infty$:
$\E[\mathcal{F}_{\mathrm{EVE}}] \to \max_k \mathcal{F}_k$.
\end{theorem}

\begin{proof}
Since $w_k \geq 0$ and $\sum_k w_k = 1$, the expectation is a
convex combination, bounded below by the minimum and above by the
maximum.
The softmax concentrates mass on the largest $\mathcal{F}_k$ as
$\beta_{\mathrm{sel}} \to \infty$ (standard property of the softmax
function).
\end{proof}

\begin{corollary}[EVE is at least as good as Adam]
Since $\mathbf{d}_1$ is the Adam direction and EVE's expected fitness
satisfies $\E[\mathcal{F}_{\mathrm{EVE}}] \geq \min_k \mathcal{F}_k$,
EVE never underperforms the worst offspring.
When $\beta_{\mathrm{sel}}$ is sufficiently large, EVE selects the
best-performing direction, which may be Adam's or may be a better
alternative.
In the worst case where Adam is always the best strategy, EVE converges
to selecting $\mathbf{d}_1$ with $w_1 \to 1$, recovering Adam.
\end{corollary}

\subsection{Convergence at $K = 1$}

\begin{proposition}
At $K = 1$, EVE is exactly AdamW (\cref{prop:k1}).
Therefore all convergence guarantees established for
Adam~\cite{kingma2015adam} and
AdamW~\cite{loshchilov2019decoupled}---including those of
\cite{reddi2018convergence} with the $\beta_2$ correction---transfer
directly to EVE at $K = 1$.
\end{proposition}

\subsection{Selection Diversity}

\begin{proposition}[Non-degenerate selection]
\label{prop:diversity}
For finite $\beta_{\mathrm{sel}}$ and $K \geq 2$, every offspring
receives positive weight: $w_k > 0\;\forall\,k$.
No offspring direction is ever fully excluded from the update.
\end{proposition}

\begin{proof}
For finite $\beta_{\mathrm{sel}}$,
$\exp(\beta_{\mathrm{sel}}\,\mathcal{F}_k) > 0\;\forall\,k$,
so $w_k > 0$ by definition of softmax.
\end{proof}

This is the ``soft selection'' property: no hard elitism, no individual
direction is ever completely discarded.
Even a poorly performing offspring contributes a small exploratory
component, maintaining search diversity.

\subsection{Dimension Recovery}

\begin{proposition}[Interpolated complementary offspring activates neglected dimensions]
\label{prop:recovery}
Consider a dimension $d$ where Adam has stalled:
$|m_{t,d}| \approx 0$, $v_{t,d}$ large, $s_{t,d} \approx 0$.
Then:
\begin{itemize}[leftmargin=1.6em, itemsep=2pt]
  \item The baseline offspring contributes $|d_{1,d}| \approx 0$ on
    dimension $d$ (stalled).
  \item The slow momentum offspring contributes $|d_{2,d}| \approx 0$
    (also stalled, since slow momentum also averages toward zero on
    sign-oscillating dimensions, albeit more slowly).
  \item The interpolated complementary offspring contributes
    $|d_{3,d}| \approx (1 - \alpha) \cdot |g_{t,d}| \cdot 1.0
    / (\sqrt{v_{t,d}} + \varepsilon)$---half the fresh gradient signal,
    with full strength activation, providing a noise-dampened but
    non-zero push on dimension $d$.
\end{itemize}
If $\mathbf{d}_3$ achieves positive fitness on the probe, it receives
positive selection weight, breaking the stall on dimension $d$.
\end{proposition}

\subsection{Sharpness Sensitivity of the Virtual SAM Offspring}

\begin{proposition}[Virtual SAM detects sharpness via probe]
\label{prop:sam}
Let $\lambda_{\max}$ denote the largest eigenvalue of the local Hessian
$\nabla^2 \Loss(\theta_t)$.
By Taylor expansion:
\begin{equation*}
  \Loss(\theta_t + \eta\,\mathbf{d}_4) \approx
  \Loss(\theta_t + \eta\,\mathbf{d}_1)
  + \eta\gamma\, g_t^\top \sign(g_t)
  + \frac{(\eta\gamma)^2}{2}\, \sign(g_t)^\top
    \nabla^2 \Loss \cdot \sign(g_t)
\end{equation*}
The first-order perturbation term $\eta\gamma\,\|g_t\|_1$ adds a
constant offset.
The second-order term is
$\frac{(\eta\gamma)^2}{2}\,\sign(g_t)^\top H\,\sign(g_t)$,
which is large when the Hessian has large eigenvalues aligned with
$\sign(g_t)$---precisely in sharp regions.
Therefore in sharp regions $\mathbf{d}_4$ incurs higher loss than
$\mathbf{d}_1$, receives lower fitness, and is downweighted by
soft selection.
\end{proposition}


% ══════════════════════════════════════════════════════════════════════════
\section{Proposed Implementation}
\label{sec:implementation}

\subsection{Full Algorithm}

\begin{algorithm}[H]
\caption{EVE: Evolutionary Virtual Exploration.}
\label{alg:eve}
\begin{algorithmic}[1]
\Require Learning rate $\eta$, decay rates $\beta_1, \beta_2,
  \beta_{1,\mathrm{slow}}$,
  brood size $K$, weight decay $\lambda$, strength decay $\gamma_s$,
  entropy target ratio $\rho$, interpolation $\alpha$,
  SAM scale $\gamma$
\State Initialize $\theta_0$, $m_0 \leftarrow 0$, $v_0 \leftarrow 0$,
  $m_{0,\mathrm{slow}} \leftarrow 0$,
  $s_0 \leftarrow 0.5 \cdot \mathbf{1}$,
  $\beta_{\mathrm{sel}} \leftarrow 1.0$
\For{$t = 1, 2, \ldots, T$}
  \State Sample batch $B_t$; compute $g_t = \nabla_\theta
    \Loss(\theta_t; B_t)$
    \hfill {\color{graytext}\textit{standard backward pass}}
  \State $m_t \leftarrow \beta_1 m_{t-1} + (1 - \beta_1)\,g_t$;
    \quad $v_t \leftarrow \beta_2 v_{t-1} + (1 - \beta_2)\,g_t^2$
  \State $m_{t,\mathrm{slow}} \leftarrow
    \beta_{1,\mathrm{slow}} m_{t-1,\mathrm{slow}}
    + (1 - \beta_{1,\mathrm{slow}})\,g_t$
    \hfill {\color{graytext}\textit{slow momentum buffer}}
  \State Bias-correct:
    $\hat{m}_t$, $\hat{v}_t$, $\hat{m}_{t,\mathrm{slow}}$
  \If{$K = 1$} \hfill {\color{graytext}\textit{exact AdamW, skip probe}}
    \State $\theta_{t+1} \leftarrow \theta_t
      - \eta\,\hat{m}_t / (\sqrt{\hat{v}_t} + \varepsilon)
      - \eta\lambda\,\theta_t$
  \Else
    \State Construct offspring $\mathbf{d}_1, \ldots, \mathbf{d}_K$
      \hfill (\cref{eq:offspring_baseline}--\cref{eq:offspring_sam})
    \State Save $\hat{\theta} \leftarrow \theta_t$
      \hfill {\color{graytext}\textit{one parameter backup}}
    \State $L_{\mathrm{base}} \leftarrow \Loss(\theta_t;\, B_t)$
      \hfill {\color{graytext}\textit{baseline forward, no backward}}
    \For{$k = 1, \ldots, K$}
      \State $\theta \leftarrow \theta_t + \eta\,\mathbf{d}_k$
        \hfill {\color{graytext}\textit{perturb in-place}}
      \State $L_k \leftarrow \Loss(\theta;\, B_t)$
        \hfill {\color{graytext}\textit{probe forward, no backward}}
      \State $\theta \leftarrow \hat{\theta}$
        \hfill {\color{graytext}\textit{restore}}
    \EndFor
    \State $\mathcal{F}_k \leftarrow L_{\mathrm{base}} - L_k$
      \quad $\forall\,k$
    \State $w_k \leftarrow \softmax(\beta_{\mathrm{sel}} \cdot
      \mathcal{F}_k)$ \quad $\forall\,k$
      \hfill (\cref{eq:selection_weights})
    \State $\theta_{t+1} \leftarrow \theta_t
      + \eta \sum_{k=1}^K w_k\,\mathbf{d}_k
      - \eta\lambda\,\theta_t$
      \hfill (\cref{eq:eve_update})
  \EndIf
  \State Update strength signal $s_{t+1}$
    \hfill (\cref{eq:strength})
  \State Adapt $\beta_{\mathrm{sel}}$ via entropy targeting
    \hfill (\cref{eq:temperature_adapt})
\EndFor
\State \Return $\theta_T$
\end{algorithmic}
\end{algorithm}

\subsection{Memory Overhead}

EVE adds buffers beyond AdamW as follows:

\medskip
\begin{center}
\small
\begin{tabular}{@{}lcc@{}}
  \toprule
  Buffer & AdamW & EVE \\
  \midrule
  Parameters $\theta$ & $D$ & $D$ \\
  First moment $m$ & $D$ & $D$ \\
  Second moment $v$ & $D$ & $D$ \\
  Slow momentum $m_{\mathrm{slow}}$ & --- & $D$ \\
  Strength signal $s$ & --- & $D$ \\
  Offspring matrix $\mathbf{D}_K$ & --- & $KD$ (temporary) \\
  Parameter backup $\hat{\theta}$ & --- & $D$ (temporary) \\
  \midrule
  Persistent & $3D$ & $5D$ \\
  Peak (during probe) & $3D$ & $5D + (K+1)D$ \\
  \bottomrule
\end{tabular}
\end{center}

\medskip
The slow momentum buffer and strength signal add two persistent
$D$-sized allocations.
The offspring matrix and parameter backup are allocated only during the
probe and released immediately after selection.
For $K = 1$ there is no probe; the persistent overhead is
$+67\%$ ($5D$ vs.\ $3D$).

\subsection{Computational Cost Summary}

\begin{center}
\small
\begin{tabular}{@{}lccc@{}}
  \toprule
  & AdamW & EVE ($K = 1$) & EVE ($K > 1$) \\
  \midrule
  Forward passes & 1 & 1 & $2 + K$ \\
  Backward passes & 1 & 1 & 1 \\
  Per-param ops & $\mathcal{O}(D)$ & $\mathcal{O}(D)$
    & $\mathcal{O}(KD)$ \\
  Wall-clock overhead & --- & 0\% & $\sim\!\frac{1+K}{3}\times 100\%$ \\
  \bottomrule
\end{tabular}
\end{center}

\notebox{%
\textbf{Where the overhead comes from.}
The $2 + K$ forward passes are: one for the gradient (shared with
AdamW), one baseline forward at $\theta_t$ to establish $L_{\mathrm{base}}$,
and $K$ probe forwards---one per offspring.
For $K = 4$ this adds $\sim$167\% overhead relative to AdamW.
The overhead scales linearly with $K$ and with the model's forward-pass
cost relative to the full step (fwd + bwd).
}


% ══════════════════════════════════════════════════════════════════════════
\section{Hyperparameter Reference}
\label{sec:hyperparams}

\begin{center}
\small
\begin{tabular}{@{}llp{7.2cm}@{}}
  \toprule
  Symbol & Default & Description \\
  \midrule
  \multicolumn{3}{@{}l}{\textit{Standard (inherited from AdamW)}} \\[3pt]
  $\eta$ & task-dependent & Learning rate \\
  $\beta_1$ & $0.9$ & First moment decay \\
  $\beta_2$ & $0.999$ & Second moment decay \\
  $\varepsilon$ & $10^{-8}$ & Numerical stabilizer \\
  $\lambda$ & $0.01$ & Decoupled weight decay \\[6pt]
  \multicolumn{3}{@{}l}{\textit{EVE-specific}} \\[3pt]
  $K$ & $4$ & Brood size (number of offspring directions).
    $K = 1$ recovers AdamW exactly. \\
  $\beta_{1,\mathrm{slow}}$ & $0.999$ &
    Slow momentum decay for $\mathbf{d}_2$.
    Higher values produce smoother trajectory tracking. \\
  $\alpha$ & $0.5$ &
    Momentum interpolation for $\mathbf{d}_3$.
    $\alpha = 1$: pure momentum (no anti-stall);
    $\alpha = 0$: pure gradient (maximum anti-stall,
    higher noise). \\
  $\gamma$ & $0.01 \cdot \eta$ &
    SAM perturbation scale for $\mathbf{d}_4$.
    Scales with learning rate to maintain relative perturbation size. \\
  $\gamma_s$ & $0.99$ & Strength signal decay rate \\
  $\rho$ & $0.5$ & Target entropy ratio for
    $\beta_{\mathrm{sel}}$ adaptation
    ($\rho = 0$: pure exploitation;
     $\rho = 1$: pure exploration) \\
  $\alpha_\beta$ & $0.01$ & Selection temperature adaptation rate \\
  $\beta_{\min}, \beta_{\max}$ & $0.1,\; 100$ &
    Selection temperature clamp \\
  \bottomrule
\end{tabular}
\end{center}

\warnbox{%
\textbf{Minimal tuning surface.}
The primary hyperparameter requiring user attention is $K$.
The offspring-specific parameters ($\beta_{1,\mathrm{slow}}$, $\alpha$,
$\gamma$) are designed with robust defaults derived from established
values in the literature: $\beta_{1,\mathrm{slow}} = 0.999$ mirrors
standard slow EMA practice, $\alpha = 0.5$ is the neutral midpoint,
and $\gamma = 0.01 \cdot \eta$ follows SAM's perturbation scaling
convention.
Whether these require task-specific tuning is an empirical question to
be investigated.
$K = 4$ is recommended as the default; $K = 1$ for baselines;
$K = 8$ for high-risk optimization (noisy gradients, sharp landscapes).
}


% ══════════════════════════════════════════════════════════════════════════
\section{Analysis of Expected Behavior}
\label{sec:analysis}

The following analyses describe \emph{expected} behavior based on the
mathematical structure of EVE's components.
Empirical validation is required to confirm these predictions.

\subsection{Expected Offspring Dominance Profile}

The offspring set is designed to produce a balanced dominance
profile where no offspring is dead (0\% dominance) and no offspring
is overloaded ($>$50\% dominance).
A healthy selection profile would show each offspring dominant
in the regime it was designed for:

\begin{center}
\small
\begin{tabular}{@{}llp{8cm}@{}}
  \toprule
  Offspring & Expected Regime & Notes \\
  \midrule
  $\mathbf{d}_1$ (Baseline)
    & Well-conditioned landscapes
    & Should carry the plurality of dominance on smooth tasks \\
  $\mathbf{d}_2$ (Slow Momentum)
    & Oscillatory trajectories
    & Active when fast momentum oscillates between valley walls \\
  $\mathbf{d}_3$ (Complementary)
    & Stalled dimensions
    & Active on tasks with high sign-oscillation in gradients \\
  $\mathbf{d}_4$ (Virtual SAM)
    & Wide minima
    & Survives selection in flat regions; rejected in sharp ones \\
  \bottomrule
\end{tabular}
\end{center}

\subsection{Smooth Convex Landscapes}

In smooth, well-conditioned settings, $\mathbf{d}_1$ and
$\mathbf{d}_2$ are expected to be near-identical (since fast and slow
momentum converge to the same direction when the gradient is stable).
The probe would assign roughly equal weight to both, and the virtual
SAM offspring would survive selection (flat region, perturbation harmless).
Only the complementary offspring would receive low weight (no stalled
dimensions to recover).

\subsection{Oscillatory Landscapes (Valley Walls)}

When the optimization trajectory oscillates between steep valley walls,
the fast momentum $\hat{m}_t$ oscillates in sign on the
oscillation-prone dimensions.
The slow momentum $\hat{m}_{t,\mathrm{slow}}$ averages over
$\sim$10$\times$ more steps ($\beta_{1,\mathrm{slow}} = 0.999$ vs.\
$\beta_1 = 0.9$), filtering out the oscillation.
The probe would detect that $\mathbf{d}_2$ reduces loss more
consistently than $\mathbf{d}_1$ in this regime, shifting selection
weight toward trajectory smoothing---the same effect Lookahead achieves
through its slow-weight interpolation, but triggered adaptively by
the probe rather than on a fixed schedule.

\subsection{Sharp Minima}

In a sharp minimum, the virtual SAM offspring $\mathbf{d}_4 =
\mathbf{d}_1 + \gamma \cdot \sign(g_t)$ probes a perturbation in
the neighborhood of the Adam step.
If the region is sharp, the perturbation increases loss:
$\mathcal{F}_4 < \mathcal{F}_1$, and $\mathbf{d}_4$ is downweighted.
The combined update shifts away from $\mathbf{d}_4$'s direction,
implicitly biasing the optimizer away from the sharp minimum.
This provides a softer, lower-cost form of sharpness awareness
compared to SAM's full second forward--backward pass.

\subsection{Hyperparameter Robustness}

At suboptimal learning rates, the offspring are expected to respond
differently:
\begin{itemize}[leftmargin=1.6em, itemsep=3pt]
  \item $\eta$ too large: The slow momentum offspring $\mathbf{d}_2$
    provides natural dampening (heavier smoothing absorbs the
    overshoot). Selection would shift toward $\mathbf{d}_2$.
  \item $\eta$ too small: The complementary offspring $\mathbf{d}_3$
    injects fresh gradient signal on stalled dimensions, partially
    compensating for insufficient step size.
  \item $\eta$ correct: Selection would balance all four offspring
    based on local geometry, with $\mathbf{d}_1$ and $\mathbf{d}_2$
    sharing the majority of weight.
\end{itemize}

This implicit adaptation is hypothesized to widen the basin of
acceptable learning rates.
Whether EVE at a suboptimal $\eta$ can match the performance of Adam at
a tuned $\eta$ is a key empirical question for the proposed experiments.


% ══════════════════════════════════════════════════════════════════════════
\section{Evolutionary Interpretation}
\label{sec:evolutionary}

EVE's design philosophy rests on a distinction between
\emph{population evolution} and \emph{strategy evolution}.
Traditional evolutionary optimizers (ES, CMA-ES, neuroevolution)
maintain a population of candidate parameter vectors.
EVE instead evolves the \emph{update strategy} applied to a single
parameter vector: the offspring are not candidate solutions but
candidate \emph{behaviors}.
This reframing eliminates the $\mathcal{O}(N)$ memory and
$\mathcal{O}(N^2)$ interaction costs of population-based methods
while preserving the core evolutionary advantage---competing strategies
are evaluated on empirical fitness and selection pressure drives
adaptation.
The strength signal and temperature adaptation play the role of
mutation and genetic drift: they ensure that the selection process
does not prematurely converge to a single strategy.

The brood is designed so that each offspring occupies a distinct
ecological niche---no two offspring serve the same function, and
no offspring is structurally non-viable.
The pilot study (\cref{sec:pilot}) demonstrated the cost of
violating this principle: a non-viable offspring (raw gradient)
forces the remaining viable offspring to bear disproportionate
selection load, collapsing the brood's effective diversity from
$K = 4$ to $K \approx 2$.

The softmax selection over offspring strategies is also a discrete
approximation to the replicator equation from evolutionary game
theory~\cite{taylor1978evolutionary}:
$\dot{x}_k = x_k\,(\mathcal{F}_k - \bar{\mathcal{F}})$,
where $x_k$ is the relative weight of strategy $k$ and
$\bar{\mathcal{F}}$ is mean fitness.
EVE's selection weights evolve analogously across steps.


% ══════════════════════════════════════════════════════════════════════════
\section{Research Directions and Open Questions}
\label{sec:discussion}

\subsubsection{Sub-Batch Probing}

The current implementation evaluates every offspring on the full
training batch $B_t$, adding $K$ full forward passes per step
($\sim$167\% overhead for $K = 4$).
An important near-term direction is \emph{sub-batch probing}: each
offspring is evaluated on a disjoint fraction $B_k \subset B_t$ of
size $|B_t|/K$, so the total probe cost is exactly one additional
full-batch-equivalent forward pass regardless of $K$.
This reduces overhead to $\sim$33\%, independent of brood size.
The key question is whether a sub-batch of size $|B_t|/K$ provides
sufficient signal quality for reliable offspring ranking---particularly
at small batch sizes.

\subsubsection{AMSGrad as Infrastructure Upgrade}

Adam's second moment $v_t$ can decrease over time as old large gradients
decay out of the EMA, causing the effective learning rate on some
dimensions to swing erratically~\cite{reddi2018convergence}.
AMSGrad addresses this by replacing $\hat{v}_t$ with the running
maximum $\hat{v}_t^{\max} = \max(\hat{v}_{t-1}^{\max}, \hat{v}_t)$,
guaranteeing monotonically non-increasing learning rates per dimension.
Since the denominator $\sqrt{\hat{v}_t} + \varepsilon$ is shared across
all offspring, switching to AMSGrad would benefit every offspring
simultaneously.
However, the current pilot evidence (CIFAR-10 / ResNet-18) does not
isolate non-monotonic $v_t$ as the primary cause of dimension neglect.
If future experiments on harder tasks (ImageNet, language modeling)
show persistent overloading of $\mathbf{d}_3$, AMSGrad should be
investigated as a shared infrastructure change.

\subsubsection{Adaptive Momentum Interpolation}

The interpolation coefficient $\alpha = 0.5$ in $\mathbf{d}_3$ is
a neutral default.
Task-specific or schedule-based adaptation of $\alpha$ (e.g.,
higher $\alpha$ early in training for stability, lower $\alpha$ late
for fine-grained recovery) is a natural extension.
However, any adaptive scheme for $\alpha$ must avoid coupling with
the strength signal to prevent feedback loops---$\alpha$ should be
driven by training progress metrics (epoch fraction, loss plateau
detection) rather than by $s_t$ itself.

\subsubsection{Optimal Offspring Design}

The four canonical offspring cover: stable exploitation ($\mathbf{d}_1$),
trajectory smoothing ($\mathbf{d}_2$), dimension recovery
($\mathbf{d}_3$), and sharpness awareness ($\mathbf{d}_4$).
For $K > 4$, additional offspring could include:
sign-momentum directions (analogous to Lion~\cite{chen2024lion}),
layer-aware offspring that apply different strategies to attention
heads vs.\ feed-forward layers, or stochastic perturbations of
$\mathbf{d}_1$ for explicit exploration.

\subsubsection{Interaction with Learning Rate Schedules}

EVE's adaptive selection partially subsumes the role of learning rate
warmup and annealing.
Whether EVE can fully replace learning rate schedules, or whether it
should be composed with them, is an empirical question.
Early experiments should test EVE with constant $\eta$ against
Adam with cosine annealing.

\subsubsection{Second-Order Offspring}

The current offspring use only first-order information ($g_t$, $m_t$,
$v_t$, $s_t$).
An offspring direction based on approximate Hessian-vector products
(e.g., via finite differences: $(g(\theta + \varepsilon p) - g(\theta))
/ \varepsilon$ for a probe direction $p$) would add curvature
information beyond $v_t$'s diagonal approximation.
This requires an additional backward pass per offspring, changing the
cost calculus; a single Hessian offspring at $K = 5$ may be worthwhile.

\subsubsection{Distributed Training}

In data-parallel distributed training, the gradient $g_t$ is
all-reduced across workers.
The probe evaluation can be distributed similarly: each worker
evaluates $K / W$ offspring on its local sub-batch shard, and
fitness values are all-reduced.
This would maintain EVE's constant overhead across worker counts.

\subsubsection{Formal Convergence for $K > 1$}

While $K = 1$ inherits Adam's convergence guarantees, the full
$K > 1$ dynamics---where the update direction changes adaptively
at each step---require a separate convergence analysis.
The key challenge is that the effective update direction is a function
of the loss surface (through the probe), introducing a non-trivial
dependence between the update and the stochastic gradient.
We conjecture that convergence holds under standard smoothness and
bounded-variance assumptions, with the selection mechanism acting
as an adaptive preconditioner, but a formal proof is left to future
work.


% ══════════════════════════════════════════════════════════════════════════
\section{Expected Contributions}
\label{sec:conclusion}

This proposal presents EVE, an optimizer that embeds evolutionary selection
pressure directly in the parameter update rule.
By constructing $K$ offspring directions---competing update strategies
derived from existing moment buffers---and selecting among them via
probe-based fitness evaluation on the training batch, EVE transforms
Adam's rigid single-formula update into an adaptive multi-strategy system.

The framework is parameterized by a single integer $K$:
at $K = 1$, EVE is exactly AdamW with zero overhead;
at $K > 1$, it evaluates $K$ strategies via forward-only probes on the
training batch ($\sim$167\% additional compute for $K = 4$), with
sub-batch probing as a clear path to reducing this cost.
The strength signal provides evolutionary memory across steps,
enabling the complementary offspring to target dimensions where Adam's
momentum has stalled.

The offspring direction set is informed by pilot evaluation, which
established two firm design constraints: momentum smoothing is
non-negotiable for training stability, and each offspring must occupy
a distinct strategic niche to prevent selection collapse.
The four canonical offspring---baseline AdamW, slow-momentum trajectory
smoother, interpolated complementary for dimension recovery, and
virtual SAM for sharpness awareness---satisfy these constraints while
maintaining EVE's zero-memory and single-backward-pass efficiency
guarantees.

EVE's design philosophy is that evolution is not about populations---it
is about \emph{selection pressure applied to competing strategies}.
A population of one is sufficient, provided the strategies are diverse
and the selection is informed by real loss on real data.
The offspring ensure diversity (by construction), and the softmax
ensures softness (no strategy is ever fully discarded).

The proposed experimental program will determine whether this theoretical
promise translates into practical gains across architectures, tasks, and
scales.


% ── References ─────────────────────────────────────────────────────────────
\begin{thebibliography}{25}

\bibitem{kingma2015adam}
Kingma, D.~P.\ \& Ba, J.\ (2015).
Adam: A method for stochastic optimization.
In \textit{Proceedings of ICLR 2015}.

\bibitem{loshchilov2019decoupled}
Loshchilov, I.\ \& Hutter, F.\ (2019).
Decoupled weight decay regularization.
In \textit{Proceedings of ICLR 2019}.

\bibitem{foret2021sharpnessaware}
Foret, P., Kleiner, A., Mobahi, H., \& Neyshabur, B.\ (2021).
Sharpness-aware minimization for efficiently improving generalization.
In \textit{Proceedings of ICLR 2021}.

\bibitem{reddi2018convergence}
Reddi, S.~J., Kale, S., \& Kumar, S.\ (2018).
On the convergence of Adam and beyond.
In \textit{Proceedings of ICLR 2018}.

\bibitem{liu2020radam}
Liu, L., Jiang, H., He, P., Chen, W., Liu, X., Gao, J., \& Han, J.\ (2020).
On the variance of the adaptive learning rate and beyond.
In \textit{Proceedings of ICLR 2020}.

\bibitem{shazeer2018adafactor}
Shazeer, N.\ \& Stern, M.\ (2018).
Adafactor: Adaptive learning rates with sublinear memory cost.
In \textit{Proceedings of ICML 2018}.

\bibitem{chen2024lion}
Chen, X., Liang, C., Huang, D., Real, E., Wang, K., Liu, Y.,
Pham, H., Dong, X., Luber, T., Cho, M., Bengio, S., \& Le, Q.~V.\ (2024).
Symbolic discovery of optimization algorithms.
In \textit{Proceedings of NeurIPS 2023}.

\bibitem{defazio2023dadaptation}
Defazio, A.\ \& Mishchenko, K.\ (2023).
Learning-rate-free learning by D-adaptation.
In \textit{Proceedings of ICML 2023}.

\bibitem{mishchenko2024prodigy}
Mishchenko, K.\ \& Defazio, A.\ (2024).
Prodigy: An expeditiously adaptive parameter-free learner.
In \textit{Proceedings of ICML 2024}.

\bibitem{defazio2024schedulefree}
Defazio, A., Yang, X., Mehta, H., Mishchenko, K., Khaled, A.,
\& Cutkosky, A.\ (2024).
The road less scheduled.
\textit{arXiv:2405.15682}.

\bibitem{zhang2019lookahead}
Zhang, M.~R., Lucas, J., Hinton, G., \& Ba, J.\ (2019).
Lookahead optimizer: $k$ steps forward, 1 step back.
In \textit{Proceedings of NeurIPS 2019}.

\bibitem{izmailov2018averaging}
Izmailov, P., Podoprikhin, D., Garipov, T., Vetrov, D., \&
Wilson, A.~G.\ (2018).
Averaging weights leads to wider optima and better generalization.
In \textit{Proceedings of UAI 2018}.

\bibitem{mansilla2021domain}
Mansilla, L., Echeveste, R., Milone, D.~H., \& Ferrante, E.\ (2021).
Domain generalization via gradient surgery.
In \textit{Proceedings of ICCV 2021}.

\bibitem{salimans2017evolution}
Salimans, T., Ho, J., Chen, X., Sidor, S., \& Sutskever, I.\ (2017).
Evolution strategies as a scalable alternative to reinforcement learning.
\textit{arXiv:1703.03864}.

\bibitem{hansen2001cma}
Hansen, N.\ \& Ostermeier, A.\ (2001).
Completely derandomized self-adaptation in evolution strategies.
\textit{Evolutionary Computation}, 9(2), 159--195.

\bibitem{taylor1978evolutionary}
Taylor, P.~D.\ \& Jonker, L.~B.\ (1978).
Evolutionary stable strategies and game dynamics.
\textit{Mathematical Biosciences}, 40(1--2), 145--156.

\bibitem{zhuang2022surrogate}
Zhuang, J., Gong, B., Yuan, L., Cui, Y., Adam, H., Dvornek, N.,
Tatikonda, S., Duncan, J., \& Liu, T.\ (2022).
Surrogate gap minimization improves sharpness-aware training.
In \textit{Proceedings of ICLR 2022}.

\bibitem{liu2022towards}
Liu, Y., Mai, S., Chen, X., Hsieh, C.-J., \& You, Y.\ (2022).
Towards efficient and scalable sharpness-aware minimization.
In \textit{Proceedings of CVPR 2022}.

\bibitem{andrychowicz2016l2l}
Andrychowicz, M., Denil, M., Gomez, S., Hoffman, M.~W., Pfau, D.,
Schaul, T., Shillingford, B., \& de~Freitas, N.\ (2016).
Learning to learn by gradient descent by gradient descent.
In \textit{Proceedings of NeurIPS 2016}.

\bibitem{metz2019understanding}
Metz, L., Maheswaranathan, N., Nixon, J., Freeman, D., \&
Sohl-Dickstein, J.\ (2019).
Understanding and correcting pathologies in the training of learned
optimizers.
In \textit{Proceedings of ICML 2019}.

\bibitem{metz2022velo}
Metz, L., Harrison, J., Freeman, C.~D., Merchant, A., Beyer, L.,
Bradbury, J., Agrawal, N., Poole, B., Mordatch, I., Roberts, A.,
\& Sohl-Dickstein, J.\ (2022).
VeLO: Training versatile learned optimizers by scaling up.
\textit{arXiv:2211.09760}.

\bibitem{kaddour2023velo}
Kaddour, J., Zhu, X., Liu, B., \& Hospedales, T.\ (2023).
Is scaling learned optimizers worth it? Evaluating the value of
VeLO's 4000 TPU months.
\textit{arXiv:2310.18191}.

\bibitem{bello2017neural}
Bello, I., Zoph, B., Vasudevan, V., \& Le, Q.~V.\ (2017).
Neural optimizer search with reinforcement learning.
In \textit{Proceedings of ICML 2017}.

\end{thebibliography}

\end{document}
